% A LaTeX template for MSc Thesis submissions to 
% Politecnico di Milano (PoliMi) - School of Industrial and Information Engineering
%
% S. Bonetti, A. Gruttadauria, G. Mescolini, A. Zingaro
% e-mail: template-tesi-ingind@polimi.it
%
% Last Revision: October 2021
%
% Copyright 2021 Politecnico di Milano, Italy. NC-BY

\documentclass{Configuration_Files/PoliMi3i_thesis}

%------------------------------------------------------------------------------
%	REQUIRED PACKAGES AND  CONFIGURATIONS
%------------------------------------------------------------------------------

% CONFIGURATIONS
\usepackage{parskip} % For paragraph layout
\usepackage{setspace} % For using single or double spacing
\usepackage{emptypage} % To insert empty pages
\usepackage{multicol} % To write in multiple columns (executive summary)
\setlength\columnsep{15pt} % Column separation in executive summary
\setlength\parindent{0pt} % Indentation
\raggedbottom  

% PACKAGES FOR TITLES
\usepackage{titlesec}
% \titlespacing{\section}{left spacing}{before spacing}{after spacing}
\titlespacing{\section}{0pt}{3.3ex}{2ex}
\titlespacing{\subsection}{0pt}{3.3ex}{1.65ex}
\titlespacing{\subsubsection}{0pt}{3.3ex}{1ex}
\usepackage{color}

% PACKAGES FOR LANGUAGE AND FONT
\usepackage[english]{babel} % The document is in English  
\usepackage[utf8]{inputenc} % UTF8 encoding
\usepackage[T1]{fontenc} % Font encoding
\usepackage[11pt]{moresize} % Big fonts

% PACKAGES FOR IMAGES
\usepackage{graphicx}
\usepackage{transparent} % Enables transparent images
\usepackage{eso-pic} % For the background picture on the title page
\usepackage{subfig} % Numbered and caption subfigures using \subfloat.
\usepackage{tikz} % A package for high-quality hand-made figures.
\usetikzlibrary{}
\graphicspath{{./Images/}} % Directory of the images
\usepackage{caption} % Coloured captions
\usepackage{xcolor} % Coloured captions
\usepackage{amsthm,thmtools,xcolor} % Coloured "Theorem"
\usepackage{float}

% STANDARD MATH PACKAGES
\usepackage{amsmath}
\usepackage{amsthm}
\usepackage{amssymb}
\usepackage{amsfonts}
\usepackage{bm}
\usepackage[overload]{empheq} % For braced-style systems of equations.
\usepackage{fix-cm} % To override original LaTeX restrictions on sizes

% PACKAGES FOR TABLES
\usepackage{tabularx}
\usepackage{longtable} % Tables that can span several pages
\usepackage{colortbl}

% PACKAGES FOR ALGORITHMS (PSEUDO-CODE)
\usepackage{algorithm}
\usepackage{algorithmic}

% PACKAGES FOR REFERENCES & BIBLIOGRAPHY
\usepackage[colorlinks=true,linkcolor=black,anchorcolor=black,citecolor=black,filecolor=black,menucolor=black,runcolor=black,urlcolor=black]{hyperref} % Adds clickable links at references
\usepackage{cleveref}
\usepackage[square, numbers, sort&compress]{natbib} % Square brackets, citing references with numbers, citations sorted by appearance in the text and compressed
\bibliographystyle{abbrvnat} % You may use a different style adapted to your field

% OTHER PACKAGES
\usepackage{pdfpages} % To include a pdf file
\usepackage{afterpage}
\usepackage{lipsum} % DUMMY PACKAGE
\usepackage{fancyhdr} % For the headers
\fancyhf{}

% Input of configuration file. Do not change config.tex file unless you really know what you are doing. 
% Define blue color typical of polimi
\definecolor{bluepoli}{cmyk}{0.4,0.1,0,0.4}

% Custom theorem environments
\declaretheoremstyle[
  headfont=\color{bluepoli}\normalfont\bfseries,
  bodyfont=\color{black}\normalfont\itshape,
]{colored}

% Set-up caption colors
\captionsetup[figure]{labelfont={color=bluepoli}} % Set colour of the captions
\captionsetup[table]{labelfont={color=bluepoli}} % Set colour of the captions
\captionsetup[algorithm]{labelfont={color=bluepoli}} % Set colour of the captions

\theoremstyle{colored}
\newtheorem{theorem}{Theorem}[chapter]
\newtheorem{proposition}{Proposition}[chapter]

% Enhances the features of the standard "table" and "tabular" environments.
\newcommand\T{\rule{0pt}{2.6ex}}
\newcommand\B{\rule[-1.2ex]{0pt}{0pt}}

% Pseudo-code algorithm descriptions.
\newcounter{algsubstate}
\renewcommand{\thealgsubstate}{\alph{algsubstate}}
\newenvironment{algsubstates}
  {\setcounter{algsubstate}{0}%
   \renewcommand{\STATE}{%
     \stepcounter{algsubstate}%
     \Statex {\small\thealgsubstate:}\space}}
  {}

% New font size
\newcommand\numfontsize{\@setfontsize\Huge{200}{60}}

% Title format: chapter
\titleformat{\chapter}[hang]{
\fontsize{50}{20}\selectfont\bfseries\filright}{\textcolor{bluepoli} \thechapter\hsp\hspace{2mm}\textcolor{bluepoli}{|   }\hsp}{0pt}{\huge\bfseries \textcolor{bluepoli}
}

% Title format: section
\titleformat{\section}
{\color{bluepoli}\normalfont\Large\bfseries}
{\color{bluepoli}\thesection.}{1em}{}

% Title format: subsection
\titleformat{\subsection}
{\color{bluepoli}\normalfont\large\bfseries}
{\color{bluepoli}\thesubsection.}{1em}{}

% Title format: subsubsection
\titleformat{\subsubsection}
{\color{bluepoli}\normalfont\large\bfseries}
{\color{bluepoli}\thesubsubsection.}{1em}{}

% Shortening for setting no horizontal-spacing
\newcommand{\hsp}{\hspace{0pt}}

\makeatletter
% Renewcommand: cleardoublepage including the background pic
\renewcommand*\cleardoublepage{%
  \clearpage\if@twoside\ifodd\c@page\else
  \null
  \AddToShipoutPicture*{\BackgroundPic}
  \thispagestyle{empty}%
  \newpage
  \if@twocolumn\hbox{}\newpage\fi\fi\fi}
\makeatother

%For correctly numbering algorithms
\numberwithin{algorithm}{chapter}

%----------------------------------------------------------------------------
%	NEW COMMANDS DEFINED
%----------------------------------------------------------------------------


%----------------------------------------------------------------------------
%	ADD YOUR PACKAGES (be careful of package interaction)
%----------------------------------------------------------------------------

%----------------------------------------------------------------------------
%	ADD YOUR DEFINITIONS AND COMMANDS (be careful of existing commands)
%----------------------------------------------------------------------------

\newcommand{\easyfly}{EasyFly}
\newcommand{\HRI}{Human-Robot Interactions}
\newcommand{\hri}{Human-Robot interactions}
\newcommand{\HDI}{Human-Drone Interactions}
\newcommand{\hdi}{Human-Drone interactions}
\newcommand{\Dronearena}{Drone Arena Challenge}

%----------------------------------------------------------------------------
%	BEGIN OF YOUR DOCUMENT
%----------------------------------------------------------------------------

\begin{document}

\fancypagestyle{plain}{%
\fancyhf{} % Clear all header and footer fields
\fancyhead[RO,RE]{\thepage} %RO=right odd, RE=right even
\renewcommand{\headrulewidth}{0pt}
\renewcommand{\footrulewidth}{0pt}}

%----------------------------------------------------------------------------
%	TITLE PAGE
%----------------------------------------------------------------------------

\pagestyle{empty} % No page numbers
\frontmatter % Use roman page numbering style (i, ii, iii, iv...) for the preamble pages

\puttitle{
	title=Programming Environment for \\Human Drone Interaction: EasyFly, % Title of the thesis
	name=Matteo Plona, % Author Name and Surname
	course=Computer Science Engineering, % Study Programme (in Italian)
	ID  = 952967,  % Student ID number (numero di matricola)
	advisor= Prof. Luca Mottola, % Supervisor name
	%coadvisor={Name Surname}, % Co-Supervisor name, remove this line if there is none
	academicyear={2023-24},  % Academic Year
} % These info will be put into your Title page 

%----------------------------------------------------------------------------
%	PREAMBLE PAGES: ABSTRACT (inglese e italiano), EXECUTIVE SUMMARY
%----------------------------------------------------------------------------
\startpreamble
\setcounter{page}{1} % Set page counter to 1

% ABSTRACT IN ENGLISH
\chapter*{Abstract} 
TODO



%Here goes the Abstract in English of your thesis followed by a list of keywords.
%The Abstract is a concise summary of the content of the thesis (single page of text)
%and a guide to the most important contributions included in your thesis.
%The Abstract is the very last thing you write.
%It should be a self-contained text and should be clear to someone who hasn't (yet) read the whole manuscript.
%The Abstract should contain the answers to the main scientific questions that have been addressed in your thesis.
%It needs to summarize the adopted motivations and the adopted methodological approach as well as the findings of your work and their relevance and impact.
%The Abstract is the part appearing in the record of your thesis inside POLITesi,
%the Digital Archive of PhD and Master Theses (Laurea Magistrale) of Politecnico di Milano.
%The Abstract will be followed by a list of four to six keywords.
%Keywords are a tool to help indexers and search engines to find relevant documents.
%To be relevant and effective, keywords must be chosen carefully.
%They should represent the content of your work and be specific to your field or sub-field.
%Keywords may be a single word or two to four words.
\textbf{Keywords:} here, the keywords, of your thesis % Keywords

% ABSTRACT IN ITALIAN
\chapter*{Abstract in lingua italiana}
TODO



%Qui va l'Abstract in lingua italiana della tesi seguito dalla lista di parole chiave.
\textbf{Parole chiave:} qui, vanno, le parole chiave, della tesi % Keywords (italian)

%----------------------------------------------------------------------------
%	LIST OF CONTENTS/FIGURES/TABLES/SYMBOLS
%----------------------------------------------------------------------------

% TABLE OF CONTENTS
\thispagestyle{empty}
\tableofcontents % Table of contents 
\thispagestyle{empty}
\cleardoublepage

%-------------------------------------------------------------------------
%	THESIS MAIN TEXT
%-------------------------------------------------------------------------

\addtocontents{toc}{\vspace{2em}} % Add a gap in the Contents, for aesthetics
\mainmatter % Begin numeric (1,2,3...) page numbering

% --------------------------------------------------------------------------
% CHAPTERS 
% --------------------------------------------------------------------------
\chapter{Introduction}
\label{ch:intro}
In recent years, the rapid advancement of unmanned aerial vehicles, commonly known as drones,
has revolutionized various industries and opened up new possibilities for applications ranging
from aerial photography and package delivery to search and rescue operations. As drones become
increasingly integrated into our daily lives, it is crucial to explore and understand the dynamics
of their interaction with humans.

In modern drone applications, the human figure is marginal with respect to the drone.
The former usually plays the supervisor role, where the main task is to control the activities and check that every operation is carried out correctly. 
At the same time, the latter performs almost all requested tasks automatically.
This significant discrepancy undoubtedly leads to a decrease in interactions between the two.

This thesis will describe EasyFly, an accessible and high-level programming environment for drone applications.
The purpose is to provide a programming environment to allow users with different levels of expertise to experiment with drones. 
Differently from modern applications, our environment allows humans and drones to work closely together,
making EasyFly a perfect tool for conducting research in the field of human-drone interactions.


\section{The Problem: Programming Human-Drone Interactions}\label{sec:the_problem}
Human-Drone Interaction (HDI) is a branch of the more general field of Human-Robot Interaction (HRI), and it can be defined as 
`the study field focused on understanding, designing, and evaluating drone systems for use by or with human user~\cite{tezza2019hdi}`.
While the field of HRI offers valuable insights, the distinctive ability of drones to move freely in three-dimensional
space, along with their unique shapes, sets HDI as a distinct and independent area of research.

The rapid technological progress in this field has made drones increasingly efficient and autonomous in performing various
tasks. While on the one hand, these advancements enable the integration of drones into everyday life, streamlining processes
and reducing the time required for specific activities, on the other hand, modern drone applications do not represent the
ideal prototype for conducting research in the field of HDI.

The first limitation of modern drone applications in this discipline's study is that these applications are designed
to operate in large environments with minimal human presence. An example can be represented by autonomous delivery
drones~\cite{singireddy2018primeAir} or 3D mapping applications~\cite{nex3Dmapping} where the interactions with the human are purposely
reduced to the bare minimum.

In these types of applications, interactions between humans and drones are often limited to simple tasks, such as package
delivery or crop monitoring. These interactions are often repetitive and lack the diversity and complexity required in
research on HDI. Since tasks and interactions are repetitive, users are usually trained to interact with the drone
in a specific way. This training can reduce the variability in HDI, making it less suitable for research purposes.

Programming HDI is usually the field's most complex and expensive task. 
Usually, researchers in this field are unfamiliar with programming at low-level embedded systems like drones.
For this reason, a specialized team of researchers is typically required to program a custom drone application that addresses the complexity of the interactions requested. 

Moreover, the implementation phase is usually the bottleneck of the entire process; every small change to the interaction model can result in days or weeks for implementing the desired behavior.
In fact, one of the most significant challenges during the programming of HDI is the testing phase.
Modern drones are usually fragile and expensive, while tests are likely to fail. 
This phase usually introduces a high consumption of resources, both in terms of costs and time needed for repairing the entire setup before another attempt.


\section{The Solution: EasyFly}\label{sec:the_solution}
To overcome the issues related to the programming of HDI, we introduce EasyFly, a programming environment for drone applications that addresses all the research needs in the field of HDI.

To better understand the contribution of EasyFly to this research field, let us take a step back and describe the needs of researchers.

The ideal prototype of a programming environment for researchers studying HDI should address and solve all the problems related to developing drone applications used for research. 
In particular, this prototype should ultimately reduce the time and costs associated with the development phase and increase the research's effectiveness.

The first characteristic of the ideal prototype is to reduce the level of expertise needed to develop the desired drone application. 
This feature allows researchers to easily implement all the required functionality without requiring a specialized drone programmer team. 
Moreover, this feature would open the doors to a brand-new type of research where the users interact with the drone programmed by themselves.
EasyFly provides this feature by offering a set of simple operations, which indeed allows the creation of very complex behaviors.
In addition to this, EasyFly allows programming in a descriptive fashion; in this way, programs would be self-explaining and easily interpreted by anyone.

As in any other field, research on HDI should be dynamic. 
In other words, to gather all the possible insights from an interaction, the drone application must rapidly change and adapt to the situation.
If the application development cycle is too long, there is the risk of losing many possible opportunities to experiment with possible alternative solutions. 
The ideal prototype should provide the maximum flexibility in adapting to many possible situations.
For this reason, EasyFly has adopted a modular approach for both the hardware and the software components.
At any moment, a module (either software or hardware) can be attached or detached to compose the best configuration needed at that specific moment.

Last but not least, facilitating the interaction between the human and the drone should be the primary goal.
The ideal prototype should be the first promoter of the interaction. 
It should offer the best possible condition to allow the two entities to establish an interaction safely and free from any potential bias determined by the programming environment.
In other words, it should use a typology of drones that allow close contact.

To implement EasyFly, we have targeted a specific typology of unmanned aerial vehicles: nano-drones.
As the name suggests, nano-drones are simply drones with very small size and weight. Their small size makes them the best choice to facilitate human-drone interaction.

In the first place, nano-drones are less intimidating and intrusive than larger drones, making it easier for
researchers to observe how individuals react and interact with them. They also allow minimal disturbance in the
observation environment, making them perfect for avoiding any possible noise in the experiment.
Given that the drone and the human are supposed to work closely together, any possible malfunction can cause an unexpected drone crash, 
especially while experimenting with new solutions. It is easy to deduce that the smaller the drone is, the safer the interaction.
The last observation is that nano-drones are usually less expensive than bigger ones, allowing researchers
to experiment with interactions with multiple drones without affecting their budget.


\section{The Benchmark: Drone Arena Challenge}\label{sec:the_benchmark}
In the HDI domain, the research's core part, especially from the computer science perspective,
is the experimental phase. During this phase, researchers put their ideas and prototypes to test
and assess the practicality of innovative interaction models.

For a programming environment like EasyFly, testing and evaluating in a real research scenario in HDI is essential. 
The testing in real scenarios can help detect possible weaknesses in the programming environment, allowing for fine-tuning the model.

To best evaluate our EasyFly programming environment, we had the possibility to participate in the Digital Futures Drone Arena project~\cite{dronearena}.
This project allowed us to perform an in-depth analysis of the impact of using EasyFly while developing human-drone interactions. 

Drone Arena is an interdisciplinary research project that aims to create a technological and conceptual platform for interdisciplinary investigations of drones at the intersection of mobile robotics,
 autonomous systems, machine learning, and Human-Computer Interaction.
The project has three inter-related objectives:
\begin{enumerate}
    \item   The constructions of a novel aerial drone testbed that is geared towards application-level
            functionality rather than low-level control mechanisms.
    \item   The organization of two challenges where multiple teams are involved and tasked to realize
            functionality that pushes the state of the art.
    \item   Conducting empirical investigation of Human Drone Interactions in the Drone Arena.
            This includes observation, interviews, and micro-sociological video analysis to inform future
            competitions in the drone arena and to develop insights from the movement-based explorations of drone piloting.
\end{enumerate}

In these settings, our EasyFly programming environment is focused on the first two objectives.
In particular, our programming environment was one of the core parts of the novel aerial drone testbed used for the entire duration of the project. 
For the second objective of the project, we had the possibility to actively participate in the first of the two challenges organized for the Drone Arena project. 
During this challenge, we conducted a complete and in-depth evaluation of the impact of using EasyFly; 
in particular, we compared our programming environment with a simpler and lower-level one.


\section{State of the Art}\label{sec:intro_soa}
The field of HDI is an active and evolving area of research with a focus on improving the ways in which humans
and drones interact. It is a multidisciplinary field with two main research areas: technological and sociological.
Each area focuses on distinct aspects of the interaction between humans and drones~\cite{hri2009davidMaya}.

In the technological area, at the intersection of computer science, mobile robotics, autonomous systems, and machine learning,
the key focus is developing and improving the hardware and software components of drones and their interfaces~\cite{kolling2012towards, giusti2012distributed}.
The main goal of this area is to enhance the capabilities and functionality of drones to make them more user-friendly and efficient~\cite{cauchard2015droneAndMe}.

In the sociological area, which includes disciplines like social engineering, art, ethics, and political science,
the core objective is to understand how the presence and use of drones impact society, individuals, and communities~\cite{eriksson2020ethicsInMovement, anderson2012accidentally}.

Especially in research focused on the sociological area, where researchers usually have less familiarity with programming tools, 
the prototyping phase is the most complex and time-consuming. 
EasyFly tries to overcome all the issues related to this phase by creating a simple and flexible programming environment for drone applications.
Moreover, it tries to offer a new perspective in the investigation of human-drone interactions where the user plays the role of the programmer.

\section{Overview}\label{sec:intro_overview}

TODO at the end


\chapter{Chapter one}
\label{ch:chapter_one}%
% The \label{...}% enables to remove the small indentation that is generated, always leave the % symbol.

In this chapter additional useful information are reported.

\section{Sections and subsections}
\label{sec:section_name}
Chapters are typically subdivided into sections and subsections, and, optionally,
subsubsections, paragraphs and subparagraphs.
All can have a title, but only sections and subsections are numbered.
A new section is created by the command
\begin{verbatim}
\section{Title of the section}
\end{verbatim}
The numbering can be turned off by using \verb|\section*{}|.
\\
A new subsection is created by the command
\begin{verbatim}
\subsection{Title of the subsection}
\end{verbatim}
and, similarly, the numbering can be turned off by adding an asterisk as follows 
\begin{verbatim}
\subsection*{}
\end{verbatim}

\section{Equations}
\label{sec:eqs}
This section gives some examples of writing mathematical equations in your thesis.

Maxwell's equations read:
\begin{subequations}
    \label{eq:maxwell}
    \begin{align}[left=\empheqlbrace]
    \nabla\cdot \bm{D} & = \rho, \label{eq:maxwell1} \\
    \nabla \times \bm{E} +  \frac{\partial \bm{B}}{\partial t} & = \bm{0}, \label{eq:maxwell2} \\
    \nabla\cdot \bm{B} & = 0, \label{eq:maxwell3} \\
    \nabla \times \bm{H} - \frac{\partial \bm{D}}{\partial t} &= \bm{J}. \label{eq:maxwell4}
    \end{align}
\end{subequations}

Equation~\eqref{eq:maxwell} is automatically labeled by \texttt{cleveref},
as well as Equation~\eqref{eq:maxwell1} and Equation~\eqref{eq:maxwell3}.
Thanks to the \verb|cleveref| package, there is no need to use \verb|\eqref|.
Remember that Equations have to be numbered only if they are referenced in the text.

Equations~\eqref{eq:maxwell_multilabels1}, \eqref{eq:maxwell_multilabels2}, \eqref{eq:maxwell_multilabels3}, and \eqref{eq:maxwell_multilabels4} show again Maxwell's equations without brace:
\begin{align}
    \nabla\cdot \bm{D} & = \rho, \label{eq:maxwell_multilabels1} \\
    \nabla \times \bm{E} +  \frac{\partial \bm{B}}{\partial t} &= \bm{0}, \label{eq:maxwell_multilabels2} \\
    \nabla\cdot \bm{B} & = 0, \label{eq:maxwell_multilabels3} \\
    \nabla \times \bm{H} - \frac{\partial \bm{D}}{\partial t} &= \bm{J} \label{eq:maxwell_multilabels4}.
\end{align}

Equation~\eqref{eq:maxwell_singlelabel} is the same as before,
but with just one label:
\begin{equation}
    \label{eq:maxwell_singlelabel}
    \left\{
    \begin{aligned}
    \nabla\cdot \bm{D} & = \rho, \\
    \nabla \times \bm{E} +  \frac{\partial \bm{B}}{\partial t} &= \bm{0},\\
    \nabla\cdot \bm{B} & = 0, \\
    \nabla \times \bm{H} - \frac{\partial \bm{D}}{\partial t} &= \bm{J}.
    \end{aligned}
    \right.
\end{equation}

\section{Figures, Tables and Algorithms}
Figures, Tables and Algorithms have to contain a Caption that describe their content, and have to be properly reffered in the text.

\subsection{Figures}
\label{subsec:figures}

For including pictures in your text you can use \texttt{TikZ} for high-quality hand-made figures,
or just include them as usual with the command
\begin{verbatim}
\includegraphics[options]{filename.xxx}
\end{verbatim}
Here xxx is the correct format, e.g. \verb|.png|, \verb|.jpg|, \verb|.eps|, \dots.

\begin{figure}[H]
    \centering
    \includegraphics[width=0.3\textwidth]{logo_polimi_scritta.eps}
    \caption{Caption of the Figure to appear in the List of Figures.}
    \label{fig:quadtree}
\end{figure}

Thanks to the \texttt{\textbackslash subfloat} command, a single figure, such as Figure~\ref{fig:quadtree},
can contain multiple sub-figures with their own caption and label, e.g. \color{black} Figure~\ref{fig:polimi_logo1} and Figure~\ref{fig:polimi_logo2}. 

\begin{figure}[H]
    \centering
    \subfloat[One PoliMi logo.\label{fig:polimi_logo1}]{
        \includegraphics[scale=0.5]{Images/logo_polimi_scritta.eps}
    }
    \quad
    \subfloat[Another one PoliMi logo.\label{fig:polimi_logo2}]{
        \includegraphics[scale=0.5]{Images/logo_polimi_scritta2.eps}
    }
    \caption[Shorter caption]{This is a very long caption you don't want to appear in the List of Figures.}
    \label{fig:quadtree2}
\end{figure}


\subsection{Tables}
\label{subsec:tables}

Within the environments \texttt{table} and  \texttt{tabular} you can create very fancy tables as the one shown in Table~\ref{table:example}.
\begin{table}[H]
    \caption*{\textbf{Title of Table (optional)}}
    \centering 
    \begin{tabular}{|p{3em} c c c |}
    \hline
    \rowcolor{bluepoli!40} % comment this line to remove the color
     & \textbf{column 1} & \textbf{column 2} & \textbf{column 3} \T\B \\
    \hline \hline
    \textbf{row 1} & 1 & 2 & 3 \T\B \\
    \textbf{row 2} & $\alpha$ & $\beta$ & $\gamma$ \T\B\\
    \textbf{row 3} & alpha & beta & gamma \B\\
    \hline
    \end{tabular}
    \\[10pt]
    \caption{Caption of the Table to appear in the List of Tables.}
    \label{table:example}
\end{table}

You can also consider to highlight selected columns or rows in order to make tables more readable.
Moreover, with the use of \texttt{table*} and the option \texttt{bp} it is possible to align them at the bottom of the page. One example is presented in Table~\ref{table:exampleC}. 

\begin{table}[H]
\centering 
    \begin{tabular}{|p{3em} | c | c | c | c | c | c|}
    \hline
%    \rowcolor{bluepoli!40}
     & \textbf{column1} & \textbf{column2} & \textbf{column3} & \textbf{column4} & \textbf{column5} & \textbf{column6} \T\B \\
    \hline \hline
    \textbf{row1} & 1 & 2 & 3 & 4 & 5 & 6 \T\B\\
    \textbf{row2} & a & b & c & d & e & f \T\B\\
    \textbf{row3} & $\alpha$ & $\beta$ & $\gamma$ & $\delta$ & $\phi$ & $\omega$ \T\B\\
    \textbf{row4} & alpha & beta & gamma & delta & phi & omega \B\\
    \hline
    \end{tabular}
    \\[10pt]
    \caption{Highlighting the columns}
    \label{table:exampleC}
\end{table}

\begin{table}[H]
\centering 
    \begin{tabular}{|p{3em} c c c c c c|}
    \hline
%    \rowcolor{bluepoli!40}
     & \textbf{column1} & \textbf{column2} & \textbf{column3} & \textbf{column4} & \textbf{column5} & \textbf{column6} \T\B \\
    \hline \hline
    \textbf{row1} & 1 & 2 & 3 & 4 & 5 & 6 \T\B\\
    \hline
    \textbf{row2} & a & b & c & d & e & f \T\B\\
    \hline
    \textbf{row3} & $\alpha$ & $\beta$ & $\gamma$ & $\delta$ & $\phi$ & $\omega$ \T\B\\
    \hline
    \textbf{row4} & alpha & beta & gamma & delta & phi & omega \B\\
    \hline
    \end{tabular}
    \\[10pt]
    \caption{Highlighting the rows}
    \label{table:exampleR}
\end{table}

\subsection{Algorithms}
\label{subsec:algorithms}

Pseudo-algorithms can be written in \LaTeX{} with the \texttt{algorithm} and \texttt{algorithmic} packages.
An example is shown in Algorithm~\ref{alg:var}.
\begin{algorithm}[H]
    \label{alg:example}
    \caption{Name of the Algorithm}
    \label{alg:var}
    \label{protocol1}
    \begin{algorithmic}[1]
    \STATE Initial instructions
    \FOR{$for-condition$}
    \STATE{Some instructions}
    \IF{$if-condition$}
    \STATE{Some other instructions}
    \ENDIF
    \ENDFOR
    \WHILE{$while-condition$}
    \STATE{Some further instructions}
    \ENDWHILE
    \STATE Final instructions
    \end{algorithmic}
\end{algorithm} 

\vspace{5mm}

\section{Theorems, propositions and lists}

\subsection{Theorems}
Theorems have to be formatted as:
\begin{theorem}
\label{a_theorem}
Write here your theorem. 
\end{theorem}
\textit{Proof.} If useful you can report here the proof.

\subsection{Propositions}
Propositions have to be formatted as:
\begin{proposition}
Write here your proposition.
\end{proposition}

\subsection{Lists}
How to  insert itemized lists:
\begin{itemize}
    \item first item;
    \item second item.
\end{itemize}
How to insert numbered lists:
\begin{enumerate}
    \item first item;
    \item second item.
\end{enumerate}

\section{Use of copyrighted material}

Each student is responsible for obtaining copyright permissions, if necessary, to include published material in the thesis.
This applies typically to third-party material published by someone else.

\section{Plagiarism}

You have to be sure to respect the rules on Copyright and avoid an involuntary plagiarism.
It is allowed to take other persons' ideas only if the author and his original work are clearly mentioned.
As stated in the Code of Ethics and Conduct, Politecnico di Milano \textit{promotes the integrity of research,
condemns manipulation and the infringement of intellectual property}, and gives opportunity to all those
who carry out research activities to have an adequate training on ethical conduct and integrity while doing research.
To be sure to respect the copyright rules, read the guides on Copyright legislation and citation styles available
at:
\begin{verbatim}
https://www.biblio.polimi.it/en/tools/courses-and-tutorials
\end{verbatim}
You can also attend the courses which are periodically organized on "Bibliographic citations and bibliography management".

\section{Bibliography and citations}
Your thesis must contain a suitable Bibliography which lists all the sources consulted on developing the work.
The list of references is placed at the end of the manuscript after the chapter containing the conclusions.
We suggest to use the BibTeX package and save the bibliographic references  in the file \verb|Thesis_bibliography.bib|.
This is indeed a database containing all the information about the references. To cite in your manuscript, use the \verb|\cite{}| command as follows:
\\
\textit{Here is how you cite bibliography entries: \cite{knuth74}, or multiple ones at once: \cite{knuth92,lamport94}}.
\\
The bibliography and list of references are generated automatically by running BibTeX \cite{bibtex}.

\chapter{Conclusions and future developments}
\label{ch:conclusions}%
A final chapter containing the main conclusions of your research/study
and possible future developments of your work have to be inserted in this chapter.

%-------------------------------------------------------------------------
%	BIBLIOGRAPHY
%-------------------------------------------------------------------------

\addtocontents{toc}{\vspace{2em}} % Add a gap in the Contents, for aesthetics
\bibliography{Thesis_bibliography} % The references information are stored in the file named "Thesis_bibliography.bib"

%-------------------------------------------------------------------------
%	APPENDICES
%-------------------------------------------------------------------------

\cleardoublepage
\addtocontents{toc}{\vspace{2em}} % Add a gap in the Contents, for aesthetics
\appendix
\chapter{Appendix A}
If you need to include an appendix to support the research in your thesis, you can place it at the end of the manuscript.
An appendix contains supplementary material (figures, tables, data, codes, mathematical proofs, surveys, \dots)
which supplement the main results contained in the previous chapters.

\chapter{Appendix B}
It may be necessary to include another appendix to better organize the presentation of supplementary material.


% LIST OF FIGURES
\listoffigures

% LIST OF TABLES
\listoftables

% LIST OF SYMBOLS
% Write out the List of Symbols in this page
\chapter*{List of Symbols} % You have to include a chapter for your list of symbols (
\begin{table}[H]
    \centering
    \begin{tabular}{lll}
        \textbf{Variable} & \textbf{Description} & \textbf{SI unit} \\\hline\\[-9px]
        $\bm{u}$ & solid displacement & m \\[2px]
        $\bm{u}_f$ & fluid displacement & m \\[2px]
    \end{tabular}
\end{table}

% ACKNOWLEDGEMENTS
\chapter*{Acknowledgements}
Here you might want to acknowledge someone.

\cleardoublepage

\end{document}
